% Page budget: one page.
\clearpage
\section{Conclusion}\label{sec:conclusion}

Strategy libraries grow because research keeps producing new policies,
modules, and recombinations worth retaining. That growth is productive, but
it is not free: storage, validation, monitoring, governance, and operational
attention are scarce resources. The economic difficulty is that a
strategy is both an operational asset and a generative asset. It can improve
the action frontier available today, and it can supply modules that create
descendants unavailable from the current frontier alone. A compression rule
that recognizes only the first role can therefore discard the source of
future invention.

This paper formulates compression as a finite combinatorial optimization
problem. Given a current library and additive resource weights, exact safe
compression minimizes retained burden while preserving both the operational
frontier and the descendant-generating closure. The projection theorem makes
these two objects the sufficient productive state and hence the correct exact
constraint: every safe-feasible library has the same productive value even
when its resource burden differs. The innovation-safe deletion certificate is
a useful local feasibility test. Rechecking that certificate after each
deletion produces a certified safe endpoint, but stepwise pruning need not
solve the global minimum-resource problem. Inclusion-wise irreducibility and
global optimality are distinct notions, so deletion order and a complete
search over feasible sublibraries can matter.

Preserving closure is economically substantive rather than a formal
refinement. Frontier-only compression may leave all immediately available
actions unchanged while severing the modules needed for a valuable
descendant. In the sharp bridge construction, dropping closure destroys all
attainable descendant value, so normalized bridge loss reaches one. Static
coverage or current payoff equivalence is consequently not a sufficient
safety criterion when retained strategies participate in later research.

Scarcity also determines which imperfectly retained library is optimal. A
hard capacity constraint expands the feasible set only when the budget can
support another relevant combination, while a resource penalty selects the
upper envelope of finitely many value--burden lines. Capacity changes and
resource-price changes therefore alter the optimal retained library at
discrete breakpoints. The resulting selections need not be nested. Moreover,
generative modules can be complements: a module may be unproductive alone
but valuable with its bridge partners. That complementarity can make the
value of capacity nonconcave even though capacity value remains monotone.

The value decomposition supplies the second step in the paper's hierarchy:
it interprets why the compression constraints matter.  Operational value
rewards actions already available, while generative value rewards future
descendants and is exposed to the delay before innovation arrives.  The
core elasticity, duration, and switching results are in Appendix~C; extended
belief-occupation and interaction results are in Online Supplements~S1--S3.
Innovation duration measures timing exposure and channel elasticities are used
only when their value denominator is positive; contribution or finite-change
measures replace them at zero and at switching thresholds.

The exact fixtures, randomized study, and financial audits illustrate these
optimization mechanisms at different evidentiary levels. In particular, the
financial comparisons between a stepwise safe library, HiGHS-reported
minimum-resource candidates, and a frontier-only library expose the
safety--compression trade-off. Held-out outcomes are retrospective opportunity
diagnostics only: they make no causal, forecasting, or alpha claim. Natural
extensions are approximate compression with an explicit safety tolerance,
endogenous information processing, continuous belief spaces, and endogenous
verification of newly generated strategies. Each would enlarge the choice set
while preserving the paper's organizing question: how much resource burden
can be removed without sacrificing the operational and inventive capacity of
the library?
\clearpage
